\documentclass[10pt,a4paper]{amsart}
\usepackage[margin=19mm]{geometry}
\usepackage{amsmath,amssymb,mathtools}
\usepackage[scr=rsfs,cal=euler]{mathalfa}
\usepackage{xfrac}
\usepackage[unicode,hidelinks,bookmarks=false]{hyperref}
\newcommand{\ie}{i.e.\ }
\newcommand{\cf}{cf.\ }
\newcommand{\resp}{respectively\ }
\newcommand{\Subsection}{\subsection}
\providecommand{\nobreakdash}{\nobreak\hbox{-}}
\setlength{\emergencystretch}{2em}
\multlinegap0pt
\usepackage{bm}
\usepackage{tikz}
\usepackage{amssymb}
\usepackage{mathtools}
\usepackage{xcolor}
\newtheorem{theorem}{Theorem}
\usepackage{cleveref}
\crefname{theorem}{Theorem}{Theorems}
\title{How iteration composition influences convergence and stability in deep learning}
\author{Benoit Dherin, Benny Avelin, Anders Karlsson, Hanna Mazzawi, Javier Gonzalvo,, Michael Munn}
\date{Source: arXiv:2502.01557v3 (CC BY 4.0)}
\begin{document}
\maketitle

\noindent\emph{Source and licence.} How iteration composition influences convergence and stability in deep learning. Version arXiv:2502.01557v3 (CC BY 4.0), distributed by arXiv under the Creative Commons Attribution 4.0 International licence, http://creativecommons.org/licenses/by/4.0/, as recorded in the arXiv licence metadata for this exact version and shown on the abstract page https://arxiv.org/abs/2502.01557v3. Full licence notice: \texttt{licenses/arxiv-2502.01557-CC-BY-4.0.txt}.\par\smallskip

\noindent\textit{Editorial scope.} Counted unit: the article's theorem lemma:contraction\_principle (source0.tex lines 707--723). The article's complete original proof of it is delivered with it (source0.tex lines 725--762, one \texttt{proof} environment of 38 lines). No mathematics is written, repaired or re-derived: the statement and the proof are the article's own lines, in the article's own notation, and no retained line is substituted or re-flowed.  No further source-local context is needed: every cross-reference of every retained line resolves inside the two delivered spans.  Nothing was omitted from any retained span, so the package carries no omission record.  No retained line contains a citation, so the wrapper carries no bibliography and no bibliography entry.  No retained line contains a figure environment, an image-inclusion command or drawing code, so no image is delivered and the package declares no asset dependency.  Editorial formatting only, in addition: an AMS \texttt{amsart} preamble that restates the article's own declarations and macros used by the retained lines, namely the environments \texttt{theorem} on separate counters, together with the standard packages those lines need.  The retained environments print in this document as Theorem 1 in order of appearance, whereas the article prints the same items with the numbering of its own full document.  The complete native source of the article as distributed in the arXiv e-print for this exact version is bundled unchanged as \texttt{sources/172665/source0.tex}.
\begin{theorem}\label{lemma:contraction_principle}
    (Backward contraction mappings principle) Let $T_{i}$ be a sequence of continuous self-maps of a complete metric space.
    Assume 
    $T_i$'s are uniform
    contractions, with a certain $k<1$ in common, and for some $\theta$ there is a constant $D$ such that,
    \begin{equation}\label{condition:boundedness}
        d(\theta,T_{i}(\theta))<D\quad \textrm{for all} \; i.
    \end{equation}
    Then for any $\theta \in \Omega$ the backward iterates 
    \[
    \theta_n = T_{1} T_{2}\cdots T_{n}(\theta)
    \]
    converge to a point $\theta^*$ as $n\rightarrow \infty$. Moreover, the convergence rate is exponential: i.e,  there is a constant $C$ depending on $\theta$ such that 
    \[
    d(\theta^*, T_{1} T_{2}\cdots T_{n}(\theta))\leq C\cdot k^{n}.
    \]
\end{theorem}

\begin{proof}
    Condition~\cref{condition:boundedness} expresses that the distance $d(\theta, T_i(\theta))$ is uniformly bounded by a constant $D$ for all $T_i's$ for a point $\theta$. Let us show that if this happens for a single point $\theta$, this happens for all points, provided we change the constant $D$. To see this, take another point $\tilde \theta$. We will compute another constant $\tilde D$ such that $d(\tilde \theta, T_i(\tilde \theta)) < \tilde D$. Namely, using the triangle inequality and the fact that $T_i$ are uniform contractions, we obtain that
    \begin{eqnarray*}
        d(\tilde \theta,T_i(\tilde \theta)) 
        & \leq & d(\tilde \theta, \theta)+d(\theta,T_i(\theta))+d(T_i( \theta),T_i (\tilde \theta)) \\
        & \leq & D+(1+k)d(\theta, \tilde \theta) = \tilde D.
    \end{eqnarray*}
    Now, we want to prove that $\theta_n$ is a Cauchy sequence, i.e, that $d(\theta_n, \theta_m)$ tends to zero for $m>n$ as $n\rightarrow \infty$. Since $\Omega$ is assumed to be complete, this will mean that the backward iterates $\theta_n$ converge toward a point $\theta^*$.  
    The idea is to bound the quantity
    \begin{eqnarray*}
    A & = &  d(\theta_n, \theta_m) \\
      & = & d(T_{1}T_{2}\cdots T_{n}(\theta),\,T_{1}T_{2}\cdots T_{m}(\theta)).
    \end{eqnarray*}
    Because of the backward order (note: the forward order would not allow that), we can apply the contraction property $n$ times (since $m> n$), yielding:
    \begin{eqnarray*}
        A & \leq & k^{n}\cdot d(\theta,\,T_{n+1}\cdots T_{m}(\theta)).
    \end{eqnarray*}
Now a simple application of the triangle inequality produces
    \begin{eqnarray*}
        A & \leq &  k^{n}\Big(d(\theta,T_{n+1}(\theta)) \\
          &      & + d(T_{n+1}(\theta),\,T_{n+1}T_{n+2}(\theta))+\cdots\\
          &      & \cdots +d(T_{n+1}\cdots T_{m-1}(\theta),\,T_{n+1}\cdots T_{m}(\theta))\Big).
    \end{eqnarray*}
At this point, we can use the condition in \cref{condition:boundedness} for the first term $d(\theta,T_{n+1}(\theta)) < D$ and in conjunction with the contraction property for the subsequent terms in the sum, yielding:
    \begin{eqnarray*}
        A  & \leq &  k^{n}\left(D+D\cdot k+\cdots + D\cdot k^{m-n+1}\right) \\
           & \leq & D\cdot k^n  \cdot \left( \frac{1 - k^{n-m+2}}{1-k}\right) \\
          & \leq & \frac{D\cdot k^{n}}{1-k},
    \end{eqnarray*}
    where we used the sum of a geometric series. Now, this yields that $A\rightarrow 0$ as $n\rightarrow\infty$, meaning that the sequence $\theta_n$
    is a Cauchy sequence and thus converges to a certain point $\theta^*$
    since $\Omega$ is complete. 
    Then, taking the limit $m\rightarrow \infty$ we obtain
    \[
    d(\theta_n,\theta^*) \leq \frac{D\cdot k^{n}}{1-k} = C\cdot k^n,
    \]
    with $C:=\frac{D}{1-k}$, which shows the exponential convergence rate.
\end{proof}

\end{document}
